\section{Annexes}\label{sec:annexes}
% Annexes (tache 10, passage resume, titre et annexes). Aucun chiffre tape: les deux
% tableaux d'hypotheses sont des corps generes par scripts/rechnung_retraite.py
% (annexe_hypotheses: data/tab_hypotheses_sources.tex depuis data/quellen.json,
% data/tab_hypotheses_modele.tex en macros de data/kennzahlen.tex); les listes de
% sources sont data/quellen_annexe.tex et data/presse_annexe.tex (quellen_tex,
% presse_notes_tex). La section "Donnees non disponibles" resume LUECKEN.md.
% La bibliographie n'est composee que si data/literatur.bib contient au moins une
% entree (\NombreReferencesAcademiques, compte par le script): sinon une phrase le dit,
% au lieu d'un titre "References" suivi d'une liste vide.

Cette annexe r\'epond \`a la question~: sur quelles hypoth\`eses et sur quelles sources reposent les
calculs de ce document, et qu'est-ce qu'il n'a pas pu \'etablir~? La
section~\ref{sec:annexe-hypotheses} r\'eunit les hypoth\`eses d\'eclar\'ees, la
section~\ref{sec:annexe-lacunes} les donn\'ees cherch\'ees sans succ\`es et ce qui les remplace, la
section~\ref{sec:annexe-perimetre} ce qui reste hors du plan, et la section~\ref{sec:annexe-sources}
la liste compl\`ete des sources cit\'ees en note.

\subsection{Hypoth\`eses d\'eclar\'ees}\label{sec:annexe-hypotheses}

Toutes les valeurs l\'egales, sociales et de march\'e du plan viennent d'un seul fichier de donn\'ees,
dont chaque entr\'ee porte sa source et sa date de consultation. Le
tableau~\ref{tab:hypotheses-sourcees} en donne les \NombreHypothesesSourcees{} entr\'ees, dont
\NombreHypothesesPrimaires{} reposent sur une source primaire (loi, administration, organisme
public)~; les autres sont des sources secondaires, signal\'ees comme telles. Chaque ligne renvoie \`a la
section qui emploie la valeur et, par un lien, \`a sa source compl\`ete de la
section~\ref{sec:annexe-sources}.

{\footnotesize
\setlength{\LTpre}{\medskipamount}\setlength{\LTpost}{\medskipamount}
\begin{longtable}{>{\bfseries\raggedright\arraybackslash}p{0.35\linewidth}
  >{\raggedright\arraybackslash}p{0.34\linewidth} >{\centering\arraybackslash}p{0.09\linewidth}
  >{\raggedright\arraybackslash}p{0.11\linewidth}}
  \caption{Hypoth\`eses sourc\'ees du plan, avec la valeur retenue, la section qui l'emploie et le type
  de source (lien vers la source compl\`ete en annexe).}\label{tab:hypotheses-sourcees}\\
  \toprule
  \rowcolor[gray]{0.9}
  \textbf{Hypoth\`ese} & \textbf{Valeur retenue} & \textbf{Section} & \textbf{Source} \\
  \midrule
  \endfirsthead
  \toprule
  \rowcolor[gray]{0.9}
  \textbf{Hypoth\`ese} & \textbf{Valeur retenue} & \textbf{Section} & \textbf{Source} \\
  \midrule
  \endhead
  \bottomrule
  \endlastfoot
  \tabellenkoerper{data/tab_hypotheses_sources}
\end{longtable}}

% Phrase liee aux donnees actuelles (a relire si quellen.json change): les sources
% secondaires sont einstiegsgehalt_brutto, gehaltssteigerung_real, etf_welt_rendite_div.
Ce tableau montre que l'essentiel du calcul repose sur des textes de loi et des publications
officielles de 2026~; les \NombreHypothesesSecondaires{} sources secondaires portent sur le salaire
d'entr\'ee, sa progression et le rendement de distribution de l'ETF monde, des valeurs dont la
section~\ref{sec:annexe-lacunes} explique l'origine. Ces valeurs sont fig\'ees \`a leur niveau de 2026 pendant tout le plan, comme le pose
l'hypoth\`ese d'ouverture du chapitre~\ref{sec:fiscalite}.

Le reste du mod\`ele repose sur des conventions propres \`a ce document, qu'aucune source ne fixe. Le
tableau~\ref{tab:conventions-modele} les r\'eunit, avec les deux \'ecarts assum\'es au plan de d\'epart~:
le rendement r\'eel du march\'e pris comme moyenne g\'eom\'etrique et non comme m\'ediane, et le
formulaire W-8BEN suppos\'e d\'epos\'e.

\begin{table}[!tb]
  \centering\small
  \caption{Conventions du mod\`ele, non sourc\'ees, avec la section qui les introduit (euros de
  2026).}\label{tab:conventions-modele}
  \begin{tabularx}{\linewidth}{>{\bfseries\raggedright\arraybackslash}p{0.3\linewidth}
    >{\raggedright\arraybackslash}X >{\centering\arraybackslash}p{0.09\linewidth}}
    \toprule
    \rowcolor[gray]{0.9}
    \textbf{Convention} & \textbf{Valeur retenue} & \textbf{Section} \\
    \midrule
    \tabellenkoerper{data/tab_hypotheses_modele}
    \bottomrule
  \end{tabularx}
\end{table}

Ce tableau montre les choix qui p\`esent le plus sur la date de ton d\'epart~: la cible, le taux
d'\'epargne, le rendement r\'eel du march\'e et la croissance du dividende de l'ETF maison. La
derni\`ere est l'hypoth\`ese la plus favorable du plan, puisqu'elle d\'epasse nettement la croissance
pass\'ee du dividende de l'indice am\'ericain (section~\ref{sec:risques-histoire}). Chaque convention
se modifie en un seul endroit du script de calcul, et le document se recalcule en entier.

\subsection{Donn\'ees non disponibles}\label{sec:annexe-lacunes}

Plusieurs donn\'ees ont \'et\'e cherch\'ees sans \^etre trouv\'ees sous la forme voulue~; le journal
d\'etaill\'e des recherches est le fichier \texttt{LUECKEN.md} du projet. Pour quatre valeurs, un
repli document\'e remplace la source id\'eale. Le texte du courrier du minist\`ere des
Finances qui fixe le Basiszins de 2026 n'a pas pu \^etre lu directement~; sa valeur de \Basiszins~\%
est confirm\'ee par deux comptes rendus ind\'ependants du m\^eme courrier. Aucune grille de salaire
propre \`a Jungheinrich ou \`a la convention collective de la m\'etallurgie n'a \'et\'e trouv\'ee pour un
ing\'enieur dipl\^om\'e en 2026~: le salaire d'entr\'ee de \GehaltBrutEntree~€ est la m\'ediane d'un
rapport de salaires priv\'e, et sa progression r\'eelle de \GehaltWachstum~\% par an est d\'eduite
d'une comparaison entre niveaux d'exp\'erience, pas d'une carri\`ere observ\'ee. La table officielle des
retenues \`a la source imputables est celle de 2025, l'\'edition de 2026 n'\'etant pas encore parue~;
les taux ne changent qu'avec les conventions fiscales. Enfin, le plus grand ETF monde est capitalisant
et ne publie aucun rendement de distribution~: le plan garde un repli de \RenditeDivEtf~\%, alors que
la part distribuante consult\'ee affichait moins (section~\ref{sec:capital-methode}).

Deux points restent \`a trancher. Pour l'inflation allemande de \InflationPicAnnee, Destatis annon\c{c}ait
d'abord \InflationPicProvisoire~\%, puis ses communiqu\'es ult\'erieurs citent \InflationPicRevisee~\%,
vraisemblablement apr\`es un changement d'ann\'ee de base de l'indice des prix~; le
chapitre~\ref{sec:risques} cite les deux et sa figure suit la s\'erie r\'evis\'ee. Par ailleurs,
plusieurs r\`egles sont employ\'ees sans entr\'ee sourc\'ee dans le fichier de donn\'ees~: les taux de
\KistAcht{} et \KistNeun~\% de l'imp\^ot d'\'Eglise, l'imputation des retenues et la d\'eduction des
cotisations dans la G\"unstigerpr\"ufung, le forfait de frais professionnels du salaire net et la table
brut-net du contr\^ole de salaire, cit\'ee en note manuelle. Elles sont signal\'ees \`a leur place et
restent \`a sourcer.

La litt\'erature acad\'emique pr\'evue au chapitre~\ref{sec:presse} n'a pas pu \^etre v\'erifi\'ee. La
base Scopus est rest\'ee inaccessible, faute de cl\'e d'acc\`es et de r\'eseau universitaire, et
l'interface publique de Semantic Scholar a refus\'e presque toutes les requ\^etes lors de deux essais
le 30~septembre~2026, les rares r\'eponses \'etant hors sujet. Aucune \'etude n'est donc cit\'ee, et
aucune ne l'est de m\'emoire~; la section~\ref{sec:annexe-references} en tient compte. C\^ot\'e presse,
Finanztest et le Handelsblatt n'ont livr\'e qu'une page de navigation ou un court chapeau derri\`ere
leur acc\`es payant, un troisi\`eme magazine \'etait ferm\'e \`a l'outil de recherche et un blog
sp\'ecialis\'e n'a pas pu \^etre joint~; deux blogs de la communaut\'e FIRE les remplacent en partie.

Le portefeuille-exemple compte \NombreTitres{} titres et non les \NombreTitresVise{} pr\'evus~: les
r\`egles de s\'election n'en laissent pas passer davantage, m\^eme rel\^ach\'ees une \`a une, et le
document garde les r\`egles plut\^ot que le nombre rond (section~\ref{sec:porte-methode}). Les donn\'ees
conserv\'ees ne disent pas quel crit\`ere a \'ecart\'e chacun des \TitresEcartesFin{} titres rejet\'es
au dernier filtre. Le nom d'une soci\'et\'e arrivait tronqu\'e des donn\'ees de march\'e~; le calcul le
coupe proprement \`a la premi\`ere virgule plut\^ot que de le compl\'eter de m\'emoire.

Certains leviers, enfin, ne sont pas mod\'elis\'es, et ce choix est d\'elib\'er\'e. La
G\"unstigerpr\"ufung vaudrait environ \GuenstigerEconomieMois~€ par mois \`a la retraite, mais le plan
garde le taux forfaitaire~: ses capitaux sont prudents sur ce point
(section~\ref{sec:fisc-guenstiger}). La rente l\'egale n'est pas d\'eduite de la cible, l'imp\^ot et
l'assurance maladie des retrait\'es ne sont pas calcul\'es, et aucun frais de gestion ou de courtage
n'est compt\'e (chapitre~\ref{sec:retraite}). Les simulations historiques ne portent que sur le
march\'e am\'ericain, et le cas d'un capital de d\'epart enti\`erement plac\'e en r\'eserve d'urgence
n'est pas recalcul\'e (section~\ref{sec:plan-premiers}).

\subsection{Hors p\'erim\`etre}\label{sec:annexe-perimetre}

Startups (Companisto), trading intraday et paris sur titres isol\'es~: hors p\'erim\`etre, aucun calcul.

\subsection{Sources}\label{sec:annexe-sources}

Les sources sont class\'ees en deux listes, chaque entr\'ee portant l'identifiant employ\'e par le
calcul, l'adresse compl\`ete et la date de consultation~; les notes de bas de page n'en affichent que
le domaine. La premi\`ere liste r\'eunit les textes de loi, les publications officielles et les donn\'ees
de march\'e, avec, pour certaines, la remarque laiss\'ee lors de la collecte~; la seconde r\'eunit les
articles de presse du chapitre~\ref{sec:presse}, consult\'es le \PresseDateConsultation.

% Coupure des URL longues des deux listes: url.sty ne coupe par defaut qu'apres "/" et
% ".", et les adresses de ces listes contiennent de longs segments avec "_" ou "-"
% (Overfull jusqu'a 360 pt constate). \AnnexeCoupuresUrl autorise la coupure apres
% chaque lettre, "_" et "-", localement dans le groupe de chaque liste (preambule
% inchange, notes de bas de page non concernees: elles n'affichent que le domaine).
\newcommand{\AnnexeCoupuresUrl}{%
  \def\UrlBreaks{\do\.\do\/\do\_\do\-\do\=\do\?\do\&%
    \do\a\do\b\do\c\do\d\do\e\do\f\do\g\do\h\do\i\do\j\do\k\do\l\do\m\do\n\do\o\do\p\do\q\do\r\do\s\do\t\do\u\do\v\do\w\do\x\do\y\do\z%
    \do\A\do\B\do\C\do\D\do\E\do\F\do\G\do\H\do\I\do\J\do\K\do\L\do\M\do\N\do\O\do\P\do\Q\do\R\do\S\do\T\do\U\do\V\do\W\do\X\do\Y\do\Z}%
  \Urlmuskip=0mu plus 1mu\relax}

\subsubsection{Sources officielles, donn\'ees et r\'ef\'erences secondaires}\label{sec:annexe-sources-officielles}

{\small\raggedright\AnnexeCoupuresUrl% genere par scripts/rechnung_retraite.py (quellen_tex) - ne pas modifier a la main
\begin{description}
\item[\texttt{abgeltungsteuer\_satz}] \phantomsection\label{src:abgeltungsteuer_satz} §32d Abs. 1 EStG. \url{https://www.gesetze-im-internet.de/estg/__32d.html}, consulté le 29.09.2026 (source primaire).
\item[\texttt{basiszins\_2026}] \phantomsection\label{src:basiszins_2026} BMF-Schreiben vom 13.01.2026, Basiszins zur Berechnung der Vorabpauschale zum 2. Januar 2026, IV C 1 - S 1980/00230/012/001. \url{https://www.bundesfinanzministerium.de/Content/DE/Downloads/BMF_Schreiben/Steuerarten/Investmentsteuer/2026-01-13-basiszins-berechnung-vorabpauschale.html}, consulté le 29.09.2026 (source secondaire). \textit{Remarque de collecte~: Zugriff auf das BMF-Schreiben selbst lieferte keinen lesbaren Zahlenwert; Wert per WebFetch aus dem Stollfuss- und Ecovis-Bericht über das gleiche Schreiben bestaetigt (Basiszins 3,20 \%). Sekundaerbestaetigung, siehe LUECKEN.md.}
\item[\texttt{durchschnittsentgelt\_2026}] \phantomsection\label{src:durchschnittsentgelt_2026} §3 Abs. 2 Sozialversicherungsrechengroessen-Verordnung 2026 (SVBezGrV 2026). \url{https://www.gesetze-im-internet.de/svbezgrv_2026/__3.html}, consulté le 29.09.2026 (source primaire).
\item[\texttt{einstiegsgehalt\_brutto}] \phantomsection\label{src:einstiegsgehalt_brutto} StepStone Gehaltsreport 2026, Medianwert Ingenieurinnen/Ingenieure mit unter einem Jahr Berufserfahrung (Datenbasis Jan. 2022 - Nov. 2025). \url{https://www.stepstone.de/e-recruiting/hr-wissen/gehalt/gehaltsreport-deutschlands-gehaelter-im-fokus/}, consulté le 29.09.2026 (source secondaire). \textit{Remarque de collecte~: Sekundaerquelle (Personaldienstleister-Report), keine unternehmensspezifische Jungheinrich/IG-Metall-Kueste-Eingruppierung gefunden, siehe LUECKEN.md}
\item[\texttt{est\_tarif\_2026}] \phantomsection\label{src:est_tarif_2026} §32a Abs. 1 EStG, Veranlagungszeitraum 2026. \url{https://www.gesetze-im-internet.de/estg/__32a.html}, consulté le 29.09.2026 (source primaire). \textit{Remarque de collecte~: Jede Zeile: [obere Grenze der Zone (EUR, null = letzte Zone), Typ, a, b, c]. Typ "null": ESt = 0. Typ "y"/"z": mit t = (x - untere Grenze der Zone)/10000, ESt = floor((a*t + b)*t + c). Typ "lin": ESt = floor(a*x + b), b bereits mit Vorzeichen (negativ) fuer die direkte Addition, quivalent zur Gesetzesschreibweise ESt = a*x - |b|. Werte gegen zwei unabhaengige Sekundaerquellen (finanz-tools.de, rechnercheck.de) kreuzgeprueft, identisch; Vorzeichen- und Typkonvention 2026-09-30 an salaire.est\_32a angeglichen (vorher progression1/progression2/linear42/linear45 mit positivem b, c=null in den linearen Zonen).}
\item[\texttt{etf\_welt\_rendite\_div}] \phantomsection\label{src:etf_welt_rendite_div} Repli explicite (Annahme), aucune valeur exploitable trouvee sur la fiche justETF du plus grand ETF MSCI World. \url{https://www.justetf.com/en/etf-profile.html?isin=IE00B4L5Y983}, consulté le 30.09.2026 (source secondaire). \textit{Remarque de collecte~: Le plus grand ETF MSCI World par encours (iShares Core MSCI World UCITS ETF USD Acc, IE00B4L5Y983, EUR 129 840 m au 31.08.2026 selon la fiche justETF) est THESAURISANT: aucun rendement de distribution n'y est affiche (verifie par WebFetch le 30.09.2026, page consultee integralement). A titre de comparaison seulement (pas retenu comme valeur, car ce n'est pas le plus grand ETF): la part distribuante iShares MSCI World UCITS ETF (Dist), IE00B0M62Q58 (https://www.justetf.com/en/etf-profile.html?isin=IE00B0M62Q58, encours EUR 8 303 m), affiche 'Current dividend yield 0,85 \%' et 'Dividend yield (1 an) 1,01 \%', tous deux sensiblement sous le repli de 1,8 \% retenu ici. Repli 0,018 conserve tel quel (decision du controleur, tache 9), ecart documente pour revision de jour. Voir LUECKEN.md.}
\item[\texttt{gehaltssteigerung\_real}] \phantomsection\label{src:gehaltssteigerung_real} Abgeleitet aus StepStone Gehaltsreport 2026 (Querschnitt nach Berufserfahrung: Einstieg 59.250 EUR vs. 25 Jahre Berufserfahrung 93.250 EUR, impliziert ca. 1,8 \%/Jahr Wachstumsrate). \url{https://www.stepstone.de/e-recruiting/hr-wissen/gehalt/gehaltsreport-deutschlands-gehaelter-im-fokus/}, consulté le 29.09.2026 (source secondaire). \textit{Remarque de collecte~: Kein direkter Destatis-Verdiensterhebungs-Zeitreihenwert zur uber die Karriere gemittelten realen Steigerung gefunden; Querschnittsdaten nach Erfahrungsjahren als Naeherung verwendet, siehe LUECKEN.md}
\item[\texttt{inflation\_2022\_de}] \phantomsection\label{src:inflation_2022_de} Statistisches Bundesamt (Destatis), Pressemitteilung Nr. 022 vom 17.01.2023, 'Inflationsrate im Jahr 2022 bei +7,9 \%'. \url{https://www.destatis.de/DE/Presse/Pressemitteilungen/2023/01/PD23_022_611.html}, consulté le 29.09.2026 (source primaire).
\item[\texttt{inflation\_destatis\_2019\_2025}] \phantomsection\label{src:inflation_destatis_2019_2025} Statistisches Bundesamt (Destatis), Pressemitteilungen annuelles 'Inflationsrate im Jahr JJJJ': PD20\_019\_611 (2019), PD21\_025\_611 (2020), PD22\_025\_611 (2021), PD25\_020\_611 (citant retrospectivement 2022 a +6,9\% et 2023 a +5,9\%), PD24\_020\_611 (2023), PD25\_020\_611 (2024), PD26\_019\_611 KORREKTUR (2025). \url{https://www.destatis.de/DE/Presse/Pressemitteilungen/2025/01/PD25_020_611.html}, consulté le 30.09.2026 (source primaire). \textit{Remarque de collecte~: Serie pour fig\_inflation\_2022 (chapitre 7, tache 9 partie B), sourcee par avance ici. ATTENTION ecart de revision: la premiere publication provisoire de janvier 2023 (PD23\_022\_611, deja citee ici sous inflation\_2022\_de) annoncait +7,9\% pour 2022; les publications retrospectives 2024/2025 (PD24\_020\_611, PD25\_020\_611) citent uniformement +6,9\% pour 2022 en comparaison des annees suivantes, probablement lie au changement de base du VPI (nouvelle base 2020=100 introduite en 2024). Serie ci-dessus prise entierement dans les publications retrospectives coherentes entre elles (2021/2022/2023/2024 cites ensemble dans PD25\_020\_611), pas melangee avec le chiffre provisoire de 2023. inflation\_2022\_de (7,9\%) laisse inchange ailleurs dans le depot (hors perimetre de cette tache); ecart signale dans LUECKEN.md pour arbitrage de jour.}
\item[\texttt{inflation\_ziel}] \phantomsection\label{src:inflation_ziel} EZB, geldpolitische Strategie, Preisstabilitaetsziel von 2 \% (HVPI). \url{https://www.ecb.europa.eu/mopo/strategy/pricestab/html/index.de.html}, consulté le 29.09.2026 (source primaire).
\item[\texttt{kv\_bbg\_monat\_2026}] \phantomsection\label{src:kv_bbg_monat_2026} GKV-Spitzenverband-Faktenblatt 'Rechengroessen und Grenzwerte im Versicherungs- und Beitragsrecht fuer das Jahr 2026', Beitragsbemessungsgrenzen Kranken- und Pflegeversicherung, S. 1. \url{https://www.gkv-spitzenverband.de/media/dokumente/presse/zahlen_und_grafiken/20260101_Faktenblatt_Rechengroessen_Beitragsrecht.pdf}, consulté le 29.09.2026 (source primaire).
\item[\texttt{kv\_mindestbemessung\_monat\_2026}] \phantomsection\label{src:kv_mindestbemessung_monat_2026} GKV-Spitzenverband-Faktenblatt 'Rechengroessen und Grenzwerte im Versicherungs- und Beitragsrecht fuer das Jahr 2026', Freiwillige Versicherung in der Krankenversicherung, Mindestbemessungsgrundlage, S. 1. \url{https://www.gkv-spitzenverband.de/media/dokumente/presse/zahlen_und_grafiken/20260101_Faktenblatt_Rechengroessen_Beitragsrecht.pdf}, consulté le 29.09.2026 (source primaire).
\item[\texttt{kv\_satz\_ermaessigt}] \phantomsection\label{src:kv_satz_ermaessigt} §243 SGB V; bestaetigt im GKV-Spitzenverband-Faktenblatt 'Rechengroessen und Grenzwerte im Versicherungs- und Beitragsrecht fuer das Jahr 2026' (1. Januar 2026), Beitragssaetze, S. 1. \url{https://www.gkv-spitzenverband.de/media/dokumente/presse/zahlen_und_grafiken/20260101_Faktenblatt_Rechengroessen_Beitragsrecht.pdf}, consulté le 29.09.2026 (source primaire). \textit{Remarque de collecte~: PDF gespeichert unter refs/gkv-spitzenverband-rechengroessen-2026.pdf}
\item[\texttt{kv\_schwellen\_dynamisierung}] \phantomsection\label{src:kv_schwellen_dynamisierung} §223 Abs. 4 SGB V: 'Die Beitragsbemessungsgrenze im Jahr 2027 entspricht der um 3 600 Euro erhoehten Jahresarbeitsentgeltgrenze nach § 6 Absatz 7 [SGB V] ... Die Bundesregierung setzt die Beitragsbemessungsgrenze in der Rechtsverordnung nach § 160 des Sechsten Buches gesondert fest.' §6 Abs. 6 Satz 2 SGB V (Jahresarbeitsentgeltgrenze, via Abs. 7 referenziert): 'Sie aendert sich zum 1. Januar eines jeden Jahres in dem Verhaeltnis, in dem die Bruttoloehne und -gehaelter je Arbeitnehmer ... im vergangenen Kalenderjahr zu den entsprechenden Bruttoloehnen und -gehaeltern im vorvergangenen Kalenderjahr stehen', d.h. eine jaehrliche Lohn-Indexierung, keine feste nominale Grenze. Die Mindestbemessungsgrundlage fuer freiwillig Versicherte ist ein Bruchteil der Bezugsgroesse (§18 SGB IV): 'das Durchschnittsentgelt der gesetzlichen Rentenversicherung im vorvergangenen Kalenderjahr, aufgerundet auf den naechsthoeheren, durch 420 teilbaren Betrag', ebenfalls jaehrlich per Sozialversicherungs-Rechengroessenverordnung festgesetzt.. \url{https://www.gesetze-im-internet.de/sgb_5/__223.html}, consulté le 30.09.2026 (source primaire). \textit{Remarque de collecte~: Weitere Fundstellen: §6 SGB V (https://www.gesetze-im-internet.de/sgb\_5/\_\_6.html, Abs. 6-7) und §18 SGB IV (https://www.gesetze-im-internet.de/sgb\_4/\_\_18.html), alle per curl am 2026-09-30 wortgetreu abgerufen. Modellkonvention (scripts/projection.py, \_saetze\_real): da Loehne mindestens so stark wie Preise wachsen und das reale Lohnwachstum im Modell separat (salaire.py) abgebildet wird, werden Mindestbemessung und BBG in scripts/projection.py KONSTANT in Euros 2026 (real) gehalten (konservativ-neutrale Wahl, keine explizite Lohnwachstumsannahme fuer die Schwellen selbst) - anders als der Sparerpauschbetrag (§20 Abs. 9 EStG), der ein fester NOMINALER Betrag ist und sich nur per Gesetzesaenderung veraendert, daher im Modell real dekumuliert (deflatiert) wird. Korrigiert eine fruehere Fehlimplementierung, die Mindestbemessung/BBG faelschlich wie den Sparerpauschbetrag deflatierte (Maximalbeitrag faellt dadurch nominal-eingefroren von ca. 1226 auf 513 EUR/Monat bis 2071, was jede Dividendenstrategie zu guenstig aussehen liess).}
\item[\texttt{kv\_teilfreistellung\_etf}] \phantomsection\label{src:kv_teilfreistellung_etf} GKV-Spitzenverband, 'Katalog von Einnahmen und deren beitragsrechtliche Bewertung nach § 240 SGB V', Stand 26.05.2026, Zeile 'Investmenterträge': '§ 16 InvStG | ja, unter Beruecksichtigung der §§ 20 und 56 Abs. 6 InvStG'. Die separate Zeile 'Kapitalvermoegen, Einkuenfte aus -' (individuelle Aktiendividenden, kein Fondsvehikel) traegt keinen solchen Verweis auf §20 InvStG., 15 von 30, Zeile 'Investmenterträge'. \url{https://www.gkv-spitzenverband.de/media/dokumente/krankenversicherung_1/grundprinzipien_1/finanzierung/beitragsbemessung/2026-05-26_Katalog_Einnahmen_beitragsrechtliche_Bewertung_240_SGB_V_BF.pdf}, consulté le 30.09.2026 (source primaire). \textit{Remarque de collecte~: Text per pdftotext -layout selbst extrahiert und gegenkontrolliert (nicht nur ueber eine KI-Zusammenfassung des Fetch-Tools uebernommen, da eine erste WebSearch-Zusammenfassung und eine erste WebFetch-Zusammenfassung sich widersprachen). Gilt nur fuer Fondsertraege (§16 InvStG, deckt sowohl Ausschuettung als auch Vorabpauschale ab); individuelle Aktiendividenden (Poche maison) erhalten keine Teilfreistellung, weder steuerlich noch in der GKV-Basis.}
\item[\texttt{kv\_zusatzbeitrag\_2026}] \phantomsection\label{src:kv_zusatzbeitrag_2026} GKV-Spitzenverband-Faktenblatt 'Rechengroessen und Grenzwerte im Versicherungs- und Beitragsrecht fuer das Jahr 2026', Beitragssaetze, Beitragssaetze, S. 1. \url{https://www.gkv-spitzenverband.de/media/dokumente/presse/zahlen_und_grafiken/20260101_Faktenblatt_Rechengroessen_Beitragsrecht.pdf}, consulté le 29.09.2026 (source primaire). \textit{Remarque de collecte~: PDF gespeichert unter refs/gkv-spitzenverband-rechengroessen-2026.pdf; vom Bundesgesundheitsministerium fuer 2026 festgelegter Schaetzerkreis-Wert}
\item[\texttt{pv\_satz\_kinderlos\_2026}] \phantomsection\label{src:pv_satz_kinderlos_2026} GKV-Spitzenverband-Faktenblatt 'Rechengroessen und Grenzwerte im Versicherungs- und Beitragsrecht fuer das Jahr 2026', Beitragssaetze (3,6 \% + 0,6 \% Zuschlag fuer Kinderlose), Beitragssaetze, S. 1. \url{https://www.gkv-spitzenverband.de/media/dokumente/presse/zahlen_und_grafiken/20260101_Faktenblatt_Rechengroessen_Beitragsrecht.pdf}, consulté le 29.09.2026 (source primaire). \textit{Remarque de collecte~: Fuer freiwillig Versicherte traegt der Versicherte den vollen Satz allein, daher hier 0,042 statt der Arbeitnehmerhaelfte}
\item[\texttt{quellensteuer}] \phantomsection\label{src:quellensteuer} Bundeszentralamt fuer Steuern (BZSt), 'Anrechenbarkeit der Quellensteuer auf Dividenden und Zinsen von Staaten, mit denen Deutschland ein Doppelbesteuerungsabkommen abgeschlossen hat', Stand 1. Januar 2025, Redaktionsschluss Juni 2025, laenderspezifische Zeilen, S. 1-25 (Belgien S.5, Daenemark S.7, Finnland S.9, Frankreich S.9, Irland S.12, Italien S.13, Kanada S.14, Niederlande S.21, Norwegen S.22, Spanien S.24, Schweiz S.23, Vereinigte Staaten S.24, Vereinigtes Koenigreich S.24). \url{https://www.bzst.de/SharedDocs/Downloads/DE/EU_OECD/anrechenbare_ausl_quellensteuer_2025.pdf}, consulté le 29.09.2026 (source primaire). \textit{Remarque de collecte~: PDF gespeichert unter refs/bzst-anrechenbare-quellensteuer-2025.pdf. Stand 1.1.2025, da 2026 zum Abrufzeitpunkt noch nicht veroeffentlicht war; DBA-Saetze aendern sich nur bei Vertragsaenderungen, siehe LUECKEN.md. DE-Zeile ist keine BZSt-Zeile, sondern aus abgeltungsteuer\_satz + soli\_satz abgeleitet.}
\item[\texttt{regelaltersgrenze}] \phantomsection\label{src:regelaltersgrenze} §35 SGB VI, Regelaltersgrenze fuer Jahrgaenge ab 1964. \url{https://www.deutsche-rentenversicherung.de/DRV/DE/Experten/Arbeitgeber-und-Steuerberater/summa-summarum/Lexikon/R/regelaltersgrenze}, consulté le 29.09.2026 (source primaire).
\item[\texttt{rentenwert\_2026}] \phantomsection\label{src:rentenwert_2026} Deutsche Rentenversicherung, Pressemitteilung 'Bundeskabinett beschliesst Rentenanpassung 2026' (29.04.2026), Anstieg von 40,79 auf 42,52 EUR. \url{https://www.deutsche-rentenversicherung.de/DRV/DE/Ueber-uns-und-Presse/Presse/Meldungen/2026/260429-rentenanpassung-2026-bundeskabinett.html}, consulté le 29.09.2026 (source primaire).
\item[\texttt{shiller\_url}] \phantomsection\label{src:shiller_url} Robert J. Shiller, Yale University, Online-Datentabelle S\&P 500 (Kurse, Dividenden, Ertraege, CPI) seit 1871. \url{http://www.econ.yale.edu/~shiller/data/ie_data.xls}, consulté le 29.09.2026 (source primaire). \textit{Remarque de collecte~: Erreichbarkeit per curl geprueft: HTTP 200, gueltige Excel-Datei (zuletzt gespeichert 2023, Autor RShiller)}
\item[\texttt{soli\_freigrenze\_2026}] \phantomsection\label{src:soli_freigrenze_2026} §3 Abs. 3 SolzG 1995 n.F. (Steuerfortentwicklungsgesetz 2024), erstmals VZ 2026. \url{https://www.buzer.de/gesetz/282/al234205-0.htm}, consulté le 29.09.2026 (source primaire).
\item[\texttt{soli\_satz}] \phantomsection\label{src:soli_satz} §4 SolzG 1995. \url{https://www.gesetze-im-internet.de/solzg_1995/__4.html}, consulté le 29.09.2026 (source primaire).
\item[\texttt{sparerpauschbetrag}] \phantomsection\label{src:sparerpauschbetrag} §20 Abs. 9 EStG, seit 2023. \url{https://www.gesetze-im-internet.de/estg/__20.html}, consulté le 29.09.2026 (source primaire). \textit{Remarque de collecte~: 2000 EUR bei Zusammenveranlagung}
\item[\texttt{sv\_arbeitnehmer}] \phantomsection\label{src:sv_arbeitnehmer} GKV-Spitzenverband-Faktenblatt 'Rechengroessen und Grenzwerte im Versicherungs- und Beitragsrecht fuer das Jahr 2026', S. 1. \url{https://www.gkv-spitzenverband.de/media/dokumente/presse/zahlen_und_grafiken/20260101_Faktenblatt_Rechengroessen_Beitragsrecht.pdf}, consulté le 29.09.2026 (source primaire). \textit{Remarque de collecte~: rv und av sind bereits die Arbeitnehmerhaelfte (halber Beitragssatz 18,6 \% bzw. 2,6 \%). kv\_allgemein und pv sind die vollen, hälftig geteilten Beitragssaetze (14,6 \% bzw. 3,6 \%); der Kinderlosenzuschlag von 0,6 \% traegt gesetzlich allein der Arbeitnehmer. kv\_zusatz (2,9 \%, voller Satz) ist eine Kopie von kv\_zusatzbeitrag\_2026, ebenfalls hälftig zu teilen (sv\_anteil in salaire.py teilt kv\_allgemein+kv\_zusatz durch 2).}
\item[\texttt{teilfreistellung\_aktienfonds}] \phantomsection\label{src:teilfreistellung_aktienfonds} §20 InvStG (Aktienfonds, natürliche Person im Privatvermögen). \url{https://www.gesetze-im-internet.de/invstg_2018/__20.html}, consulté le 29.09.2026 (source primaire).
\end{description}
\par}

\subsubsection{Presse sp\'ecialis\'ee}\label{sec:annexe-sources-presse}

{\small\raggedright\AnnexeCoupuresUrl% genere par scripts/rechnung_retraite.py (presse_notes_tex) - ne pas modifier a la main
\begin{description}
\item[\texttt{boerseonline\_fire\_portfolio}] \phantomsection\label{presse:boerseonline_fire_portfolio} Börse Online, \textit{FIRE: Aufbau eines Dividendenportfolios für finanzielle Unabhängigkeit}. \url{https://www.boerse-online.de/nachrichten/geldundvorsorge/fire-aufbau-eines-dividendenportfolios-fuer-finanzielle-unabhaengigkeit-20392175.html}, consulté le 30.09.2026.
\item[\texttt{boerseonline\_fire\_traum}] \phantomsection\label{presse:boerseonline_fire_traum} Börse Online, \textit{FIRE: Der Traum von der finanziellen Unabhängigkeit – das steckt hinter dem Investment-Megatrend}. \url{https://www.boerse-online.de/nachrichten/geldundvorsorge/fire-der-traum-von-der-finanziellen-unabhaengigkeit-das-steckt-hinter-dem-investment-megatrend-20385679.html}, consulté le 30.09.2026.
\item[\texttt{dividende\_statt\_rente}] \phantomsection\label{presse:dividende_statt_rente} Dividende statt Rente (blog FIRE), \textit{Die Dividendenstrategie}. \url{https://www.dividende-statt-rente.de/die-dividendenstrategie/}, consulté le 30.09.2026.
\item[\texttt{extraetf\_dividendenboom}] \phantomsection\label{presse:extraetf_dividendenboom} Extra-Magazin (extraETF), \textit{Dividenden-Boom 2026: Dax vor einem Rekordjahr – das müssen Anleger jetzt wissen}. \url{https://extraetf.com/de/news/etf-news/dividenden-boom-2026-warum-der-dax-vor-einem-rekordjahr-steht}, consulté le 30.09.2026.
\item[\texttt{extraetf\_fire}] \phantomsection\label{presse:extraetf_fire} Extra-Magazin (extraETF), \textit{FIRE-Bewegung}. \url{https://extraetf.com/de/wissen/fire-bewegung}, consulté le 30.09.2026.
\item[\texttt{finanznachrichten\_irrweg}] \phantomsection\label{presse:finanznachrichten_irrweg} FinanzNachrichten.de, \textit{Entnahmephase: Dividendenstrategie im Ruhestand: Ein mental bequemer Irrweg}. \url{https://www.finanznachrichten.de/nachrichten-2026-09/69580312-entnahmephase-dividendenstrategie-im-ruhestand-ein-mental-bequemer-irrweg-607.htm}, consulté le 30.09.2026.
\item[\texttt{finanztip\_geldanlage}] \phantomsection\label{presse:finanztip_geldanlage} Finanztip, \textit{Geldanlage: So legst Du Dein Geld einfach sicher an}. \url{https://www.finanztip.de/geldanlage/}, consulté le 30.09.2026.
\item[\texttt{finanztip\_thesaurierend}] \phantomsection\label{presse:finanztip_thesaurierend} Finanztip, \textit{Thesaurierende Fonds: So bequem investierst Du in wiederanlegende Fonds}. \url{https://www.finanztip.de/indexfonds-etf/thesaurierende-fonds/}, consulté le 30.09.2026.
\item[\texttt{getmad\_fire}] \phantomsection\label{presse:getmad_fire} getmad.de (blog FIRE), \textit{FIRE Bewegung}. \url{https://www.getmad.de/fire-bewegung/}, consulté le 30.09.2026.
\item[\texttt{justetf\_ausschuettend}] \phantomsection\label{presse:justetf_ausschuettend} justETF, \textit{Ausschüttende Strategien mit ETFs umsetzen}. \url{https://www.justetf.com/de/academy/ausschuettende-strategien.html}, consulté le 30.09.2026.
\item[\texttt{justetf\_entnahmestrategien}] \phantomsection\label{presse:justetf_entnahmestrategien} justETF, \textit{Entnahmestrategien im Vergleich: so geht es entspannt durch den Ruhestand}. \url{https://www.justetf.com/de/academy/entnahmestrategien.html}, consulté le 30.09.2026.
\end{description}
\par}

\ifnum\NombreReferencesAcademiques>0
  % Au moins une reference verifiee dans data/literatur.bib: bibliographie composee comme
  % sous-section non numerotee de l'annexe (option KOMA "leveldown").
  \KOMAoptions{bibliography=leveldown}
  \phantomsection\label{sec:annexe-references}
  \bibliographystyle{plain}
  \bibliography{data/literatur}
\else
  \subsection{R\'ef\'erences acad\'emiques}\label{sec:annexe-references}

  Aucune r\'ef\'erence acad\'emique n'a pu \^etre v\'erifi\'ee pour cette version du document (voir la
  section~\ref{sec:annexe-lacunes})~; cette liste reste donc vide plut\^ot que de citer une \'etude de
  m\'emoire. Elle se remplira d'elle-m\^eme d\`es que des r\'ef\'erences v\'erifi\'ees seront ajout\'ees
  au fichier bibliographique du projet.
\fi
